\addcontentsline{toc}{subsubsection}{Solution 27. IMO 2011 -- Problem 3}

\begin{dlsummary}[title={\bfseries Solution 27. IMO 2011 -- Problem 3},breakable]

\begin{proof}
Replacing $y$ by $t-x$ in the given inequality gives
\[
f(t)\leq(t-x)f(x)+f(f(x))\qquad(x,t\in\mathbb R). \tag{1}
\]

\medskip
\noindent
\textbf{Step 1. Derive a two-variable inequality.}

Set $t=f(z)$ in (1), then add the resulting inequality to the one
obtained by interchanging $x$ and $z$. The terms involving $f(f(x))$
and $f(f(z))$ cancel, yielding
\[
xf(x)+zf(z)\leq2f(x)f(z)\qquad(x,z\in\mathbb R). \tag{2}
\]
Taking $z=2f(x)$ in (2) cancels the terms involving
$f(2f(x))$ and gives
\[
xf(x)\leq0\qquad(x\in\mathbb R). \tag{3}
\]
In particular, $f(x)\leq0$ for $x>0$, while $f(x)\geq0$ for $x<0$.

\medskip
\noindent
\textbf{Step 2. Show that $f$ is nonpositive everywhere.}

Suppose $f(w)>0$ for some $w\in\mathbb R$. Taking $x=w$ in (1),
\[
f(t)\leq(t-w)f(w)+f(f(w)).
\]
Choose
\[
t<\min\!\left(0,\,w-\frac{f(f(w))}{f(w)}\right).
\]
Then $t<0$ and the right-hand side is negative, so $f(t)<0$.
This contradicts $f(t)\geq0$ for $t<0$, established in Step 1.
Therefore $f(x)\leq0$ for every real $x$.

\medskip
\noindent
\textbf{Step 3. Prove the required vanishing.}

For $x<0$, Step 1 gives $f(x)\geq0$, while Step 2 gives
$f(x)\leq0$. Hence $f(x)=0$ for all $x<0$. Using (1) with any
$x<0$ and $t<0$ now gives
\[
0=f(t)\leq(t-x)f(x)+f(f(x))=f(0),
\]
so $f(0)\geq0$. Step 2 also gives $f(0)\leq0$, and therefore
$f(0)=0$. Thus $f(x)=0$ for every $x\leq0$.
\end{proof}

\end{dlsummary}